\section{A short-segment estimate}
\label{sec:short-segment}

The traversal bound gives compactness, but by itself permits a limiting
curve to retrace part of its image.  We now obtain a second estimate: two
arcs travelling between opposite sides cannot carry incompatible signed
crossings through the same short interval.  The additional factor of the
interval length will come from a two-particle insertion.  Its matrix
entries have signs, so the geometric interpretation of its contraction
is an essential part of the argument.

Use physical coordinates $(x,y)$ in which transfer proceeds in the
$x$-direction; transfer lengths $t_-,t_+$ below count lattice columns.
For an oriented lattice path $\alpha$ and a segment $I$ of a vertical seam,
write $\iota_I(\alpha)$ for its algebraic intersection number with $I$,
using one fixed transverse sign convention.  Endpoints of $I$ are off the
graph, and endpoints of the paths being tested are off the seam.  The
direction on every double-dimer cycle is the one determined by the two
matching labels.  In particular, for a whole simple cycle $g$,
\begin{equation}
 \iota_I(g)\in\{-1,0,1\}.
 \label{eq:whole-loop-index}
\end{equation}

\begin{proposition}[Short intervals in a rectangle]
\label{prop:rectangle-short}
Let $R_\delta$ be corner-rooted odd lattice rectangles converging to a
fixed nondegenerate rectangle, with two independent uniform dimer
covers.  Fix an interior longitudinal coordinate $c$, two positive
numbers $\rho_-,\rho_+$, and a bounded segment $J$ of the transverse line
$x=c$.  The two regions $x\le c-\rho_-$ and $x\ge c+\rho_+$ are called the
deep regions; if either is disjoint from the rectangle, the events below
are empty.  The segment $J$ may extend beyond the transverse sides of the
rectangle.

For each sufficiently small $h>0$, partition $J$ into at most $C_J/h$
intervals of length at most $h$, with endpoints off the graph.  Use a seam converging to $x=c$
and the corresponding intervals when necessary.  A tested arc is a
cycle subarc, in its cycle direction, with one endpoint in each deep
region.  Let $B_{\delta,h}$ be the event that some partition interval $I$
satisfies at least one of the following conditions:
\begin{enumerate}[label=\textnormal{(\roman*)}]
 \item a tested arc $\alpha$ on a cycle $g$ has
 $\iota_I(\alpha)\ne0$ and $\iota_I(\alpha)\ne\iota_I(g)$;
 \item tested arcs on two distinct cycles both have nonzero intersection
 number with $I$.
\end{enumerate}
The arcs may otherwise travel anywhere in the rectangle.  Then
\begin{equation}
 \lim_{h\downarrow0}\limsup_{\delta\downarrow0}
       \PP(B_{\delta,h})=0.
 \label{eq:rectangle-short-limit}
\end{equation}
The conclusion is unchanged by truncating intervals at a transverse
side, or by placing their endpoints outside the rectangle.
\end{proposition}

\subsection{Marked caps and the analytic bound}

If either deep region is absent, there is nothing to prove.  Otherwise
choose two cap seams in the interior of the rectangle, $x=a$ and $x=b$, with
\[
 c-\rho_-<a<c<b<c+\rho_+.
\]
They remain a positive macroscopic distance from the middle seam and
from the longitudinal ends of the rectangle.  In the left cap an arch
is \emph{eligible} if it reaches $x\le c-\rho_-$; in the right cap it is
eligible if it reaches $x\ge c+\rho_+$.  Here an arch is a connected
nondoubled path in the cap with its two endpoints on the cap seam.
Let $E_M$ be the event that each cap has at most $M$ eligible arches,
where $M\ge1$.  Choose in each cap a fixed strictly interior line
between the deep line and the cap seam.  Every eligible arch makes two
traversals between this shallower line and the cap seam, even when the
deep line itself is a longitudinal side of the rectangle.
Proposition~\ref{prop:rectangle-crossings} therefore gives
\begin{equation}
 \lim_{M\to\infty}\limsup_{\delta\downarrow0}\PP(E_M^c)=0.
 \label{eq:cap-cutoff-tight}
\end{equation}
The same event $E_M$ will be used for all intervals of the partition.

For now assign each eligible arch $A$ a weight $\lambda(A)\in[0,1]$
determined only by its own uncolored cap diagram and the choice of $A$.
We first compute with these weights and then choose them to make the
signed contraction nonnegative.  Their dependence on one cap alone
allows the two marked cap vectors to be formed separately.

In each cap select one eligible arch and tear it as in
Lemma~\ref{lem:marked-cap}, with weight $\lambda(A)$.
Use two plus signs at the left cap and two minus signs at the right cap.
These are two \emph{sources}: their apparently different signs are due
to the opposite boundary normals of the two caps.  Sum over the
selections, and set the marked cap vector to zero when its cap has more
than $M$ eligible arches.  Denote the resulting left and right vectors
by $V_L$ and $V_R$.  In the particle convention of the preceding section
they lie in the two-replica sectors $(p+1,p-1)$ and $(p-1,p+1)$,
respectively, where $n=2p+1$ is the number of transverse sites.
In cap-vector subscripts, the seams are indexed by their column numbers,
as in Section~\ref{sec:transfer}.  Lemma~\ref{lem:marked-cap} gives
\[
 \|V_L\|\le M\|L_a\otimes L_a\|,
 \qquad \|V_R\|\le M\|R_b\otimes R_b\|.
\]

At the middle seam let $S(I)\subseteq\{1,\ldots,n\}$ be the indices of
sites in $I$.  The insertion $\mathcal O_I$ sums, over $i<j$ in $S(I)$,
the operation which removes the pair $i,j$ from the first exterior
factor and creates it in the second.  It vanishes unless both sites
have first-replica occupancy one and second-replica occupancy zero.
The resulting pair of occupancies is reversed at both sites.  Thus
\[
 \mathcal O_I:
 \Lambda^{p+1}\mathbb R^n\otimes\Lambda^{p-1}\mathbb R^n
 \longrightarrow
 \Lambda^{p-1}\mathbb R^n\otimes\Lambda^{p+1}\mathbb R^n.
\]
Its increasing-wedge sign in the site basis is
\begin{equation}
 (-1)^{\#\{\text{single occupations strictly between }i\text{ and }j\}}.
 \label{eq:insertion-parity}
\end{equation}
Indeed the two exterior factors contribute their usual deletion and
insertion signs; doubly occupied intervening indices contribute twice,
and the other nonzero contributions are precisely the single
occupations.  The operation preserves total seam-edge multiplicities
site by site, also on rows where the particle convention complements
edge occupancy.  Figure~\ref{fig:two-particle-insertion} shows the two
cap sources and the two middle-seam sinks in a same-loop contribution.

\begin{figure}[t]
\centering
\begin{tikzpicture}[x=1.35cm,y=1cm,>=Stealth,
    every node/.style={font=\small}]
  \fill[blue!4] (-4,-1.4) rectangle (-2,1.4);
  \fill[blue!4] (2,-1.4) rectangle (4,1.4);
  \draw[dashed,gray] (-2,-1.55)--(-2,1.6);
  \draw[dashed,gray] (2,-1.55)--(2,1.6);
  \draw[gray] (0,-1.55)--(0,1.6);
  \draw[line width=2pt,red!55!black] (0,-.9)--(0,.9);
  \draw[thick] (-2,.6)--(2,.6)
    .. controls (4.25,.6) and (4.25,-.6) .. (2,-.6)
    --(-2,-.6)
    .. controls (-4.25,-.6) and (-4.25,.6) .. (-2,.6);
  \fill[white] (-3.6875,0) circle (.10);
  \draw[thick] (-3.7875,0)--(-3.5875,0);
  \fill[white] (3.6875,0) circle (.10);
  \draw[thick] (3.5875,0)--(3.7875,0);
  \fill (0,.6) circle (.055);
  \fill (0,-.6) circle (.055);
  \node[above left] at (-2,.6) {$+$};
  \node[below left] at (-2,-.6) {$+$};
  \node[above right] at (2,.6) {$-$};
  \node[below right] at (2,-.6) {$-$};
  \node[above right] at (0,.6) {$j$};
  \node[below right] at (0,-.6) {$i$};
  \node[right,red!55!black] at (0,0) {$I$};
  \node at (-3,-1.15) {left cap};
  \node at (3,-1.15) {right cap};
  \node[above] at (-2,1.6) {$a$};
  \node[above] at (0,1.6) {$c$};
  \node[above] at (2,1.6) {$b$};
  \node[above] at (-3.55,.86) {source};
  \node[above] at (3.55,.86) {source};
  \draw[<->] (-1.85,-1.15)--(-.15,-1.15);
  \draw[<->] (.15,-1.15)--(1.85,-1.15);
  \node[below] at (-1,-1.15) {$t_-$};
  \node[below] at (1,-1.15) {$t_+$};
\end{tikzpicture}
\caption{A same-loop contribution to the insertion. The selected cap arches
carry two sources: the endpoint signs are \(++\) at the left seam and
\(--\) at the right seam. The two sinks \(i,j\) lie in the middle
interval \(I\). The insertion changes the permitted color assignments
while preserving the uncolored loop. The positive propagation widths
\(t_-\) and \(t_+\) control the sum over mode states.}
\label{fig:two-particle-insertion}
\end{figure}


\begin{lemma}[Sandwiched insertion]
\label{lem:sandwiched-insertion}
Let $\mathcal T_-$ and $\mathcal T_+$ be the two-replica transfers from
$a$ to $c$ and from $c$ to $b$, with column lengths $t_-$ and $t_+$.
For the fixed positive widths above, so that $t_\pm\ge\eta n$ for some
$\eta>0$ and all sufficiently small meshes,
\begin{equation}
 \bigl\|s^{-2(t_-+t_+)}
       \mathcal T_+\mathcal O_I\mathcal T_-\bigr\|
       \le C\left(\frac{|S(I)|}{n}\right)^2.
 \label{eq:sandwiched-insertion}
\end{equation}
Consequently, uniformly over the cap-local weights $\lambda(A)\in[0,1]$,
\begin{equation}
 \frac{\bigl|\langle V_R,
       \mathcal T_+\mathcal O_I\mathcal T_-V_L\rangle\bigr|}
      {Z_{\rm dd}}
 \le C_M\left(\frac{|S(I)|}{n}\right)^2.
 \label{eq:marked-contraction-bound}
\end{equation}
\end{lemma}

\begin{proof}
Use the real orthogonal singular-mode basis at the middle seam supplied
by Lemma~\ref{lem:spectrum}; it is the same for both replicas.  Write
its change-of-basis matrix as $Q$, and put
\[
 P_I=Q^T\1_{S(I)}Q.
\]
Here $\1_{S(I)}$ is the diagonal orthogonal projection onto the indicated
site coordinates.
For a removed mode pair $R=\{r,s\}$, $r<s$, and an added mode pair
$T=\{u,v\}$, $u<v$, the coefficient of pair deletion in the first factor
and pair creation in the second is
\begin{equation}\label{eq:insertion-minor}
 \sum_{\substack{i<j\\i,j\in S(I)}}
       \det Q_{\{i,j\},R}\,\det Q_{\{i,j\},T}
 =\det (P_I)_{R,T}.
\end{equation}
All pairs use increasing order; the equality is the two-by-two
Cauchy--Binet identity.  For fixed input and output occupation states,
the removed pair is their set difference in the first factor and the
added pair is their set difference in the second.  Thus a nonzero
matrix entry uses exactly one coefficient in
\eqref{eq:insertion-minor}, multiplied by the occupation-state sign.
The delocalization bound $|Q_{ir}|\le C/\sqrt n$ gives
\[
 |(P_I)_{rs}|\le C^2\frac{|S(I)|}{n},
 \qquad
 |(\mathcal O_I)_{\alpha\beta}|
       \le 2C^4\left(\frac{|S(I)|}{n}\right)^2.
\]

We must also control the sum over mode states.  For a two-replica mode
state $\mu=(J,K)$, set
\[
 d_\pm(\mu)=w_{t_\pm}(J)w_{t_\pm}(K),
\]
using the normalized single-replica weights of Lemma~\ref{lem:spectrum}.
Equation~\eqref{eq:all-state-sum}, which includes the zero mode, gives
$\sum_\mu d_\pm(\mu)\le C_\eta^2$, even when the sums run over all
particle sectors.  In the compatible singular bases the normalized
transfers are diagonal, so the sandwiched matrix has entries
$d_+(\nu)(\mathcal O_I)_{\nu\mu}d_-(\mu)$.  Consequently
\[
 \begin{split}
 \bigl\|s^{-2(t_-+t_+)}\mathcal T_+\mathcal O_I\mathcal T_-\bigr\|
 &\le\sum_{\nu,\mu}d_+(\nu)
             |(\mathcal O_I)_{\nu\mu}|d_-(\mu)\\
 &\le 2C^4\left(\frac{|S(I)|}{n}\right)^2
       \left(\sum_\nu d_+(\nu)\right)
       \left(\sum_\mu d_-(\mu)\right).
 \end{split}
\]
The operator norm is bounded by the entrywise absolute sum, and the
bounded state sums prove~\eqref{eq:sandwiched-insertion}.
The displayed bounds on $\|V_L\|,\|V_R\|$ and the partition comparison in
Lemma~\ref{lem:cap-comparison} now prove
\eqref{eq:marked-contraction-bound}.
\end{proof}

The analytic estimate supplies a factor $|S(I)|^2$ for every permitted
choice of weights.  To use it for a probability, we must choose the
weights so that the full signed contraction is nonnegative and detects
the event we want to exclude.  We now compute its contributions before
making that choice.

\subsection{The pointwise color calculation}

Fix an ordinary uncolored double-dimer configuration $\omega$.  Each
nondoubled loop has two assignments of its alternating edges to the
two matchings; doubled edges have only one.  For each selected pair of
cap arches and each sink pair $i<j$ in $S(I)$, count the modified color assignments
satisfying the two source rules and the two insertion rules.  Multiply
by~\eqref{eq:insertion-parity}, divide by the ordinary number of
assignments, and multiply by the two cap weights $\lambda$.  Sum over
the choices, and define the result to be $X^\lambda_{M,I}(\omega)$;
set it to zero on $E_M^c$.  Lemma~\ref{lem:marked-cap} gives the exact
identity
\begin{equation}\label{eq:observable-contraction}
 \EE X^\lambda_{M,I}
 =\frac{\langle V_R,\mathcal T_+\mathcal O_I\mathcal T_-V_L\rangle}
        {Z_{\rm dd}}.
\end{equation}

To compute this color ratio, cut the uncolored configuration at the
middle seam $x=c$.  A marked loop not reaching that seam has no possible
sink and contributes zero.  Otherwise the selected left cap arch lies
in a unique arch of the same loop in $\{x<c\}$.  Direct its two branches
away from the source.  Ordinary strands propagate these directions to
the middle seam, where both branches leave the left half-plane.
The induced endpoint signs are therefore $++$.  On the right the same
argument gives $--$, because both branches leave $\{x>c\}$.
An unmarked middle-seam arch has one entrance and one exit, hence
opposite endpoint signs.

Several choices of a remote eligible arch can lie on the same
middle-seam arch.  We keep them as separate marked choices, each with
its own weight.  Moving a source along that arch changes no
intersection with $I$.  Thus the color calculation below applies to
each remote choice and preserves its multiplicity when we sum.

Rank the singly occupied edges along the whole middle seam, and give
rank $r$ the alternating sign $\sigma_r=(-1)^r$.  Formula
\eqref{eq:insertion-parity} becomes
\begin{equation}
 -\sigma_i\sigma_j
 \label{eq:rank-insertion-sign}
\end{equation}
when $i,j$ now denote the two ranks.  On any one simple loop these rank
signs agree with oriented intersection signs up to a common sign.
To see why other loops do not affect this assertion, take two seam
hits of the chosen loop.  They have the same inside/outside relation
to every other disjoint loop.  Each such loop therefore contributes an
even number of hits between them.  After these even contributions are
removed, the assertion is the usual alternation of entering and
leaving a Jordan domain along a line.

\begin{lemma}[One or two marked loops]
\label{lem:color-contractions}
Fix one eligible mark in each cap, before multiplying their weights.
\begin{enumerate}[label=\textnormal{(\roman*)}]
 \item If both marks lie on one loop, let $u,v$ be the intersection
 numbers with $I$ of its two mark-to-mark arcs, both oriented from the
 left mark to the right mark.  Its summed signed color ratio is
 $uv/2$.  This is nonnegative, and is at least $1/2$ when $u,v\ne0$.
 \item If the marks lie on distinct loops $g,g'$, put
 \[
 t_g=\sum_{\text{hits of }g\text{ in }I}\sigma_r.
 \]
 Their summed signed color ratio is $-t_gt_{g'}/4$.  The number $t_g$
 is zero unless $g$ separates the endpoints of $I$, in which case it
 is $1$ or $-1$.
\end{enumerate}
\end{lemma}

\begin{proof}
One can express the color rules using signs $x_r^L,x_r^R\in\{-1,1\}$
on the two sides of each single middle-seam hit.  An unmarked arch
requires opposite signs at its ends.  A marked left arch fixes both
signs to $+1$, and a marked right arch fixes both to $-1$.  Ordinary
seam gluing requires $x_r^L=x_r^R$, whereas a selected sink fixes
$(x_r^L,x_r^R)=(+1,-1)$.  On an ordinary loop these constraints leave
one free binary choice; once a source is prescribed, propagation fixes
every remaining sign whenever the constraints are compatible.

At a source the two oriented strands point away from the mark;
at a sink they point toward it.  Continuing the assignments around a
loop shows that sources and sinks must alternate.  On a loop with two
sources, the sink pair is therefore admissible precisely when it puts
one sink on each of the two mark-to-mark arcs.  It then prescribes one
assignment instead of the two ordinary choices.  The two arcs have
opposite directions relative to any coherent orientation of the
whole loop, canceling the minus sign in
\eqref{eq:rank-insertion-sign}.  Summing the remaining products of
crossing signs yields $uv/2$.

The whole-loop intersection number, in the orientation of the first
arc, is $u-v$, so~\eqref{eq:whole-loop-index} implies $|u-v|\le1$.
Because $u,v$ are integers, their product is nonnegative, and is at
least one if neither vanishes.

For two distinct marked loops there must be one sink on each.  The
colors of each marked loop are then fixed, replacing four ordinary
choices by one.  Summing~\eqref{eq:rank-insertion-sign} gives
$-t_gt_{g'}/4$.  The rank-sign observation before the lemma identifies
$t_g$, up to a sign, with the whole-loop intersection number.  This
number vanishes exactly when the endpoints of $I$ have the same
inside/outside status with respect to $g$.
\end{proof}

For a precise parity identity, order the endpoints $A,B$ of $I$ in the
increasing seam direction, and let $d(g)$ be the number of loops strictly
enclosing $g$.  Start rank numbering at the first hit of the whole seam,
so that the region before all hits is exterior to every loop.  Then
\begin{equation}
 t_g=-(-1)^{d(g)}
 \bigl(\1_{\{B\in\operatorname{Int}(g)\}}
       -\1_{\{A\in\operatorname{Int}(g)\}}\bigr).
 \label{eq:rank-nesting-depth}
\end{equation}
Indeed at a hit entering $g$, the nesting parity just before the hit is
$d(g)$, and at a hit leaving $g$ it is $d(g)+1$.  Since
$\sigma_r=(-1)^r$, summing these signed changes telescopes to
\eqref{eq:rank-nesting-depth}, regardless of how often $g$ hits the seam.

The loops separating the endpoints of $I$ form at most two nested
chains, one around either endpoint, with their common ancestors
omitted.  Order each chain from its outermost loop inward.  The signs
$t_g$ alternate along each chain.  If both chains are nonempty, the
two outermost signs are opposite.  Successive members of either chain
have ancestor depths differing by one.  The two outermost members, when
both exist, border the same common-ancestor region and have equal
depths, but the endpoint differences in
\eqref{eq:rank-nesting-depth} have opposite signs.  Common ancestors
change both depth parities equally; other nonseparating loops do not
affect the identity.  Endpoints outside the rectangle are permitted in
this description.

\subsection{Nested arches and positivity}

Only loops separating the endpoints of $I$ contribute to the
different-loop sum, since the other loops have $t_g=0$.  For such a
loop $g$, let
\[
 U_g=\sum_{\substack{A\text{ eligible left arch}\\A\subset g}}\lambda(A),
 \qquad
 V_g=\sum_{\substack{A\text{ eligible right arch}\\A\subset g}}\lambda(A).
\]
These totals still depend on the weights to be chosen; empty sums are
zero.  The total contribution from marks on distinct
loops is, by Lemma~\ref{lem:color-contractions},
\begin{equation}
 -\frac14\sum_{g\ne g'}U_gV_{g'}t_gt_{g'}.
 \label{eq:different-loop-sum}
\end{equation}
Individual summands can be negative: two loops whose positions in a
chain differ by two have the same sign $t_g$.  We choose the weights
so that the total mark weights decrease by a factor of at least two
along either nesting chain.  The resulting alternating sums will
control the complete expression~\eqref{eq:different-loop-sum}.

The weights must still be determined within each cap.  For an eligible
arch $A$, let $b(A)$ count the eligible arches in that cap whose
endpoints strictly span those of $A$ along its seam.  Choose
\begin{equation}
 \lambda(A)=w(A):=(2M)^{-b(A)}.
 \label{eq:arch-weight}
\end{equation}
On $E_M$ every such weight belongs to $[q_M,1]$, where
\begin{equation}
 q_M=(2M)^{-(M-1)}.
 \label{eq:minimum-arch-weight}
\end{equation}
From now on $U_g,V_g$ use this choice, and write
$X_{M,I}=X^w_{M,I}$.  Every nonzero $U_g$ or $V_g$ is at least $q_M$,
because it contains at least one eligible-arch weight.
To connect the endpoint nesting chains to these
cap weights, we show that every eligible arch of an inner separating
loop is spanned by an eligible arch of the outer loop.  That additional
spanning arch supplies a factor $1/(2M)$; the cutoff of $M$ arches will
then give the required factor $1/2$ for the total weights.

\begin{lemma}[Decrease of the cap weights]
\label{lem:cap-weight-decrease}
On $E_M$, if a separating loop $g'$ lies inside another separating
loop $g$, then
\[
 U_{g'}\le\tfrac12 U_g,
 \qquad V_{g'}\le\tfrac12 V_g.
\]
In particular a zero outer weight forces all deeper weights in that
cap to be zero.
\end{lemma}

\begin{proof}
Consider an eligible arch $A'$ of $g'$ in one cap.  Since $g$ separates
the endpoints of the middle interval, it reaches the middle seam and
is not contained in the cap half-plane.  Its portions in that
half-plane are disjoint arches.  Each such arch $A$, closed by the
interval between its seam endpoints, bounds a Jordan disk $D_A$.
Inside the open cap half-plane, the indicator of the interior of $g$
is the parity sum of the indicators of these disks: the two sides
have the same boundary there and both vanish sufficiently far away.
An interior point of $A'$ therefore belongs to at least one $D_A$.
The whole interior of $A'$ stays in that disk, since it cannot cross
$A$ or the cap seam.  Thus the endpoints of $A$ strictly span those of
$A'$.

The disk $D_A$ reaches no deeper into the cap than its boundary arch.
Since $A'$ reaches the prescribed deep region, $A$ also reaches it and
is eligible.  Every eligible arch spanning $A$ also spans $A'$, and
$A$ itself supplies one additional spanning arch.  Consequently
\[
 b(A')\ge b(A)+1,
 \qquad w(A')\le\frac{w(A)}{2M}.
\]
There are at most $M$ eligible inner choices.  Bounding each containing
outer weight by the largest outer weight gives
\[
 \sum_{A'\subset g'}w(A')
 \le\tfrac12\max_{A\subset g}w(A)
 \le\tfrac12\sum_{A\subset g}w(A).
\]
This proves the assertion independently in the two caps.
\end{proof}

\begin{lemma}[Two alternating chains]
\label{lem:two-chain-positivity}
Suppose a finite family consists of at most two chains, each ordered
from outside inward.  Give its members signs $t_g\in\{-1,1\}$ which
alternate in either chain and have opposite outermost signs when
both chains are present.  Let $U_g,V_g\ge0$ decrease by a factor of at
least two at each step inward.  Then the expression
\eqref{eq:different-loop-sum} is nonnegative.  If every positive $U_g$
or $V_g$ is at least $q>0$ and there are distinct $g,g'$ with $U_gV_{g'}>0$, the
expression is at least $q^2/16$.
\end{lemma}

\begin{proof}
Within one chain write its indices as $1,\ldots,k$.  Group the terms
by their outer member to obtain
\[
 \frac14\sum_{i=1}^{k-1}
 \left[
 U_i\sum_{j>i}(-1)^{j-i+1}V_j+
 V_i\sum_{j>i}(-1)^{j-i+1}U_j
 \right].
\]
Each inner alternating sum is nonnegative and is at least half its
first term, since successive magnitudes decrease by a factor of at
least two.  If a particular product $U_iV_j$ or $V_iU_j$ with $i<j$ is
positive, the appropriate first term at $i+1$ is positive as well.
That group then contributes at least $q^2/8$.

Between the two chains the sum is
\[
 -\frac14\left[
 \left(\sum_{g\in C_1}U_gt_g\right)
 \left(\sum_{g\in C_2}V_gt_g\right)
 +
 \left(\sum_{g\in C_2}U_gt_g\right)
 \left(\sum_{g\in C_1}V_gt_g\right)
 \right].
\]
Each weighted alternating sum is zero or has its chain's outermost sign, and has
absolute value at least half its outermost weight.  The signs of the
two chains are opposite, so both displayed products give nonnegative
contributions after the leading minus sign.  A positive product
between distinct chains forces the relevant outermost weights to be
at least $q$, giving a contribution of at least $q^2/16$.
\end{proof}

We can now finish the detection argument.  Lemma~\ref{lem:color-contractions}
makes every same-loop term of $X_{M,I}$ nonnegative.
Lemmas~\ref{lem:cap-weight-decrease} and~\ref{lem:two-chain-positivity}
make its complete different-loop term nonnegative.  Thus
\begin{equation}
 X_{M,I}\ge0.
 \label{eq:positive-observable}
\end{equation}

If condition (i) of Proposition~\ref{prop:rectangle-short} occurs, each
deep endpoint of the tested arc lies in an eligible arch of its cap.
Select those two arches.  Moving the endpoints to the marks inside
these arches changes no intersection with $I$, since the arches stay
strictly on their respective sides of the middle seam.  If the tested
arc has intersection number $a$ and the whole loop has number $e$,
the two mark-to-mark sums have product $a(a-e)$.  Here $a\ne0,e$ and
$e\in\{-1,0,1\}$, so this product is a positive integer.  The weighted
same-loop term is at least $q_M^2/2$.

Suppose instead that condition (ii) occurs and no tested arc gives
condition (i).  Each of its two nonzero arc sums must then equal the
corresponding whole-loop sum.  The two loops separate the endpoints
of $I$ and each has an eligible arch in each cap.  In particular there
is a distinct pair with $U_gV_{g'}>0$.  Their full different-loop sum
is at least $q_M^2/16$.  We have proved the pointwise bound
\begin{equation}
 X_{M,I}\ge\frac{q_M^2}{16}\,
 \1_{E_M\cap\{\textnormal{condition (i) or (ii) holds on }I\}}.
 \label{eq:pointwise-detection}
\end{equation}

\subsection{Summing the short intervals}

\begin{proof}[Proof of Proposition~\ref{prop:rectangle-short}]
For a partition interval $I$, combine
\eqref{eq:marked-contraction-bound},
\eqref{eq:observable-contraction}, and
\eqref{eq:pointwise-detection} to obtain
\[
 \PP\bigl(E_M\cap\{\text{a violation on }I\}\bigr)
 \le C_M\left(\frac{|S(I)|}{n}\right)^2.
\]
The transverse size of the limiting rectangle is positive and fixed,
so an interval of length at most $h$ contains at most $C(h/\delta+1)$
sites and $n\ge c/\delta$.  Consequently
\[
 \frac{|S(I)|}{n}\le C(h+\delta).
\]
There are at most $C_J/h$ intervals.  The union bound therefore yields
\begin{equation}
 \PP(B_{\delta,h})
 \le\PP(E_M^c)+C_M\frac{(h+\delta)^2}{h}.
 \label{eq:short-union-bound}
\end{equation}
For each fixed $M$, first take the upper limit as $\delta\downarrow0$
and then let $h\downarrow0$.  The remaining bound is
$\limsup_{\delta\downarrow0}\PP(E_M^c)$, which tends to zero as
$M\to\infty$ by~\eqref{eq:cap-cutoff-tight}.  This proves the stated
order of limits.  The argument used only the sites actually contained
in an interval; endpoints outside the rectangle and intervals
truncated by a transverse side cause no change.
\end{proof}
